\chapter{两套册子}
\chapterkana{二つの名簿}
\epigraphbox{“或劳心，或劳力；劳心者治人，劳力者治于人。”}{——《孟子·滕文公上》}

\kaishou{名册上有一栏叫「手」。一页一列，一格一个数。\par
某环，某时，手若干。三个字一行，抄了几百年。\par
把这些数加起来，比全城的人还多。}

\section{两种册子}

前一章讲的《章程总目》六卷，管的是同一批人、同一座城与同一套规矩。可在同一批出土文献里，还有另一套册子，它所记录的却完全是另一批对象。

第一套是《章程总目》第四卷，即名籍卷，共一百九十二条，按一人一条登记，末尾一栏记「定日」，末枚另有一行合计：在册十一万九千。另一套则不在那四百一十二枚之中，而是单独出土的《工籍残卷》，编号 \texttt{ФВ-1968-14}。它仅存三枚，比总目那批金属片更薄也更小，字口极浅，多处笔画被锈蚀掩盖，全靠红外成像才勉强辨识。

\begin{table}[htbp]
\centering
\caption{两套册子的对照}
\begin{tabular}{@{}lll@{}}
\toprule
 & 名籍卷 & 工籍残卷 \\
\midrule
对象 & 人 & 手 \\
单位 & 一人一条 & 某环某时，报数 \\
起点 & 定日 & 不定日 \\
末栏 & 除籍 & 报 \\
合计 & 十一万九千 & 一万一千四百（某环某日） \\
\bottomrule
\end{tabular}
\tabsource{名籍卷见章程总目第四卷；工籍残卷为 \texttt{ФВ-1968-14}，三枚，第一枚轮值表、第二枚工别、第三枚上报格式。}
\end{table}

\section{缺口}

把岁计卷与这两套册子并排来看，账面上立刻露出一处无法弥合的缺口。岁计卷将全城支出分为食、燃、修、净四类，其中「食」一类的细目明写着「月给粟，人日二升」，按人头支给；其余燃、修、净三类却只记总数，不落到具体人头上。

照常理，燃料供应、设施修缮与环境清理都要人去做，做工的人总要吃饭。既然口粮按人日二升发给，十一万九千在册人口的消耗便是个定数；可真正在做燃、修、净这三类活计的人，翻遍名籍卷却一个也找不到。能找到他们的唯有《工籍残卷》。残卷第一枚是轮值表，记录某环一日内各时段的当值人数，峰值达一万一千四百，谷值也有六百；第二枚列明工别，分为农作、清扫、制药、警备四类；第三枚则是上报格式，页眉印着一行：「某环某时，手若干，报。」

一万一千四百若仅仅是某一道环某一日的峰值，那么全城六环相加，每日当值的劳力总数便已远远超出十一万九千。名籍卷上根本没有这么多人，多出来的这一大块劳力，显然不在名籍的册子里面。

\section{不在册的人}

把两套册子的体例拆开细看，便知两套册子的分别在于记录对象的性质，而不只是详略。名籍卷格式固定：先设定日一格，一人一日，绝不重复；次列迁转数行，记历任职务；末栏记除籍，理由仅限事故、逃亡、除籍三端。反观工籍残卷，三枚之上只见时刻、环号与人数三项，既无姓名，也无定日与除籍，甚至连一个「人」字都不曾出现过。

名籍上的人要分粮、分房、授籍，属于有身份的城中居民；工籍所记的则只是一笔笔劳力数目，不涉口粮与户籍，因而另立一册，仅在需要时向上报数。第三枚页眉那句「手若干，报」把这层意思挑得很明：上面过问的只是一个数字，而非具体的人。《章程总目》一千二百四十七条的边栏与正文里，「手」字共出现十七次，竟无一次用在月之民身上；凡指称月之民，一概用「民」「居者」或「在册者」三个词。

\section{造物的来历}

月之都的人口只减不增，城内的日常杂务却一项不少，账面上这个巨大缺口，只能靠另一种存在去填补。总目与工籍对此均未明言来历，笔者所能拿出来的，只是三条彼此咬合的线索。

第一条线索在工别分类。工籍残卷所列的农作、清扫、制药、警备四类，恰好一一对应城内维系运转的四套系统：供能、洁净、净秽山加工与防卫。城里需要什么，工籍上便列着什么，一项不多，一项不少。第二条线索在补充方式：第三枚上报格式旁另有一行小字，记某环当年「新添三百二十名」。一个只管报数、不管户籍的册子上出现「新添」二字，说明这批劳力的补充靠的是自身繁衍，而非外部招募。第三条线索则来自地面流传极早的一则神话：月亮上有兔子在捣药。这则传说流传的年代远早于月面材料进入档案之时，两条线并无互相抄袭的可能，却指向同一桩事实，即月面上确实存在一群会做工的兔子。\footnote{此说见《淮南子·览冥训》一系。本书只在「月亮上有兔子在干活」这一点上用它，捣的是什么、给谁捣，传说与材料各不相干。}

将这三条线索拼在一处，大致可以推出一个轮廓：城内的劳力乃是一群另行制造出来的生物，以地面上某种繁殖力极强的动物为底本，赋予其人的手与服从性，使其能生育、能劳作，却不入月之民的名籍。至于这批造物究竟是如何制成的，现有材料只字未提，只能暂时搁置。

\section{总管与报数}

工籍上报的数目究竟送往何人手中，三枚残卷未留一字，只在第三枚行末孤零零地留着一个「报」字，后面本该跟着收件处的名目，偏偏金属片在那里锈蚀残缺了。不过收件人仍可从《章程总目》中寻出端倪：名籍卷末枚有一条记着「总管一人，名不署，位不虚」，八个字点出两层意思，即这个位置常设不撤，名册上却不写究竟由谁担任，相当于一个永远占着编制、却永远隐去姓名的岗位。

这一条的刻制格式本身也露了马脚。名籍卷其余各条若无姓名，格内便直接留空；只有这一条的姓名格并不空着，而是四边封口刻成了一个方框，细看之下，像是先将原本刻好的字铲平，再补刻上一圈边框。红外成像证实铲痕深度与周边笔画一致，并非后人磨损所致。再回看工籍第三枚末尾那个「报」字，它单占一行，字号比同枚其余各字大出一号，刻痕也深出约三成。\footnote{深度用共聚焦成像量，同一枚取五点，报字那一点比其余四点平均深零点零三毫米。笔画形状与其余各字没有分别，仍是平收。}同一枚上其余各字深浅波动不及一成，只有此处用力骤增，若非临时起意，便是用以标记整段呈报规程到此为止，笔者倾向后一种读法。

至于这位总管究竟姓甚名谁，线索落在另一批外派名册上。撤离之前的一份名册记着公元前二〇五〇年曾有一人下地，职名为「地之主任」，其后括注「兼手之总管」。名册在这个名字一栏写着「嫦娥」二字，只是墨色与前后正文略有差异，像是事后补填上去的。将月面档案里的「地之主任兼总管」与地面流传了两千年的「嫦娥」接在一处固然顺当，笔者却不敢把话说死：档案只证明曾有一位地之主任兼管这批造物，尚不足以证明她便是民间故事里奔月的那一位。两边对得上一次，眼下也就只能认作对上了这一次。

复盘全章，所凭无非名籍卷、工籍残卷与一则地面神话。前两件同出一家，彼此吻合只能证明它们自己说得通；真正令笔者意外的，反倒是年代极早、来源散漫的地面神话竟能与之合拍。然而整件事里最关键的一环依然悬在半空：一座能让人不再衰老的城池，究竟是如何造出一种既能生育又听命劳作的新物种的，材料里毫无依据，全靠笔者顺着账目缺口猜出来。账无论如何算不平，可那座城的机器却实打实地转了两万年。
